\begin{proposition}[High-frequency obstruction to uniform second-order correction]
\label{prop:circle-high-frequency}
Let $M$ be the circle of radius $R>0$, equipped with uniform probability
measure, and let $A=-\Delta_M$.  Write $B_\varepsilon$ for the Sinkhorn
metric operator associated with the fixed chordal squared-distance cost,
and let
\[
 \widehat B_\varepsilon
   =B_\varepsilon+
     \frac{\varepsilon}{4R^2}B_\varepsilon A B_\varepsilon
\]
be its positive curvature correction.  For every $a>0$, set
\[
 x=\frac{a^2}{4R^2},\qquad
 c_a=\frac14\left(\frac1x-\frac1{\sinh x}\right)>0.
\]
Then, on $L^2_0(M)$,
\[
 \liminf_{\varepsilon\downarrow0}
  \frac{\|\widehat B_\varepsilon-A^{-1}\|_{L^2\to L^2}}{\varepsilon}
 \ge c_a.
\]
The same lower bound holds with $B_\varepsilon$ in place of
$\widehat B_\varepsilon$.  In particular, the corrected metric's
$O(\varepsilon^2)$ expansion on fixed smooth directions cannot be
strengthened to an $O(\varepsilon^2)$ operator-norm bound on all of
$L^2_0(M)$.
\end{proposition}

\begin{proof}
In angular coordinates, let $e_k(\theta)=\sqrt2\cos(k\theta)$ for
$k\ge1$, so $\|e_k\|_{L^2}=1$ and $Ae_k=(k^2/R^2)e_k$.
The Markov and metric eigenvalues are
\[
 p_{\varepsilon,k}=
 \frac{I_k(2R^2/\varepsilon)}{I_0(2R^2/\varepsilon)},
 \qquad
 b_{\varepsilon,k}=
 \frac{\varepsilon p_{\varepsilon,k}}
      {2(1-p_{\varepsilon,k}^2)}.
\]
Choose integers $k_\varepsilon\ge1$ with
$k_\varepsilon\sqrt\varepsilon\to a$, for example
$k_\varepsilon=\lfloor a/\sqrt\varepsilon\rfloor$ for sufficiently
small $\varepsilon$.  The integral representation of $I_k$, after
the substitution $\theta=\sqrt\varepsilon z$, gives
\[
 p_{\varepsilon,k_\varepsilon}
 =\frac{
 \displaystyle\int_{-\pi/\sqrt\varepsilon}^{\pi/\sqrt\varepsilon}
 e^{-\frac{2R^2}{\varepsilon}(1-\cos(\sqrt\varepsilon z))}
 \cos(k_\varepsilon\sqrt\varepsilon z)\,dz}
 {\displaystyle\int_{-\pi/\sqrt\varepsilon}^{\pi/\sqrt\varepsilon}
 e^{-\frac{2R^2}{\varepsilon}(1-\cos(\sqrt\varepsilon z))}\,dz}.
\]
Extend both integrands by zero outside their integration intervals.
Since $1-\cos\theta\ge2\theta^2/\pi^2$ for $|\theta|\le\pi$,
their absolute values are bounded by the integrable function
$\exp(-4R^2z^2/\pi^2)$.  The numerator integrand converges pointwise
to $e^{-R^2z^2}\cos(az)$, and the denominator integrand converges
to $e^{-R^2z^2}$.  Dominated convergence therefore yields the
triangular-array limit
\[
 p_{\varepsilon,k_\varepsilon}\longrightarrow e^{-x},
 \qquad
 \frac{b_{\varepsilon,k_\varepsilon}}{\varepsilon}
   \longrightarrow\beta_a:=\frac1{4\sinh x},
 \qquad
 \frac{R^2/k_\varepsilon^2}{\varepsilon}
   \longrightarrow\frac1{4x}.
\]
On this Fourier mode the corrected eigenvalue is
\[
 \widehat b_{\varepsilon,k}
   =b_{\varepsilon,k}
     +\frac{\varepsilon k^2}{4R^4}b_{\varepsilon,k}^2.
\]
Because $k_\varepsilon^2\varepsilon\to a^2$ and
$b_{\varepsilon,k_\varepsilon}=\varepsilon\beta_a+o(\varepsilon)$,
the added term is
$\varepsilon^2a^2\beta_a^2/(4R^4)+o(\varepsilon^2)$.
Consequently
\[
 \frac{R^2/k_\varepsilon^2-
       \widehat b_{\varepsilon,k_\varepsilon}}{\varepsilon}
 \longrightarrow\frac1{4x}-\frac1{4\sinh x}=c_a>0,
\]
where strict positivity follows from $\sinh x>x$ for $x>0$.
Testing the operator difference on the unit vector
$e_{k_\varepsilon}$ proves the claim; omitting the correction gives
the same limit.  The correction defines a bounded operator at each
fixed $\varepsilon>0$, since the Gaussian integral operator is
smoothing and its Fourier eigenvalues decay faster than every
power of $k$.
\end{proof}

\paragraph{Interpretation.}
The frequency scale $k\asymp\varepsilon^{-1/2}$ resolves the
regularisation length.  These directions depend on $\varepsilon$;
although each is smooth, their derivative norms diverge.  The lower
bound is therefore consistent with second-order convergence for
each fixed smooth direction and uniformly bounded smooth families.
Taking $a=2R$ gives the explicit positive constant
$c_a=(1-1/\sinh1)/4$.  This is a scope clarification, not a separate
novelty claim.
